%! Intercepts, Zeros, and Extrema — Unit 2, Lesson 2.4 — Group Activity
\documentclass[10pt]{article}
\usepackage{atda-article}
\usepackage{atda-boxes}
\newcommand{\TallMath}[1]{$\displaystyle #1\rule[-1.4em]{0pt}{3.2em}$}
\setlength{\workrowsep}{6pt}

\begin{document}

\pageheader{Unit 2, Lesson 2.4}{Group Activity --- Reading a Graph for Everything It Says}
% NAMESTRIP: no \namedateperiod here. The student writes their name once, on the
% cover this component is stapled behind. See references/conventions.md.

\vspace{0.15in}

\begin{scenariobox}[The Scenario]{royal}
\small
Every question today is answered by reading a graph --- no equations required, and none given.
\textbf{The rule for the whole activity: every time you write a number, say out loud whether it came
from the horizontal axis or the vertical one.} An interval of inputs and a stretch of outputs can be
the same two numbers and mean completely different things. \textbf{Everyone starts at Tier R.} When
your whole group can do Tier R without help, move on.
\end{scenariobox}

\vspace{0.10in}

\boxguard[30]
\begin{tcolorbox}[colback=white, colframe=black!40, title=\textbf{Tier R --- Remediate},
    fonttitle=\bfseries, arc=2mm, left=3mm, right=3mm, top=2mm, bottom=2mm, breakable]
\small
The graph shows $P(t)$, the number of people in the school gym, $t$ hours after it opens at
6~a.m. The gym closes when the last person leaves.

\begin{center}
\begin{tikzpicture}
\begin{axis}[
    width=10.4cm, height=4.8cm,
    xlabel={$t$ (hours after opening)}, ylabel={$P$ (people in the gym)},
    xmin=0, xmax=6, ymin=0, ymax=45,
    xtick={0,1,2,3,4,5,6}, ytick={0,5,10,15,20,25,30,35,40,45},
    grid=both, grid style={linegray!45}, axis lines=left,
    tick label style={font=\tiny}, label style={font=\scriptsize}]
  \addplot[royal, very thick] coordinates {(0,10) (1,10) (4,40) (6,0)};
\end{axis}
\end{tikzpicture}
\end{center}

\begin{enumerate}[leftmargin=1.6em, itemsep=4pt, topsep=2pt]
  \item Give the $y$-intercept as a point, and say in a sentence what it means about the gym.

\writelines{1}

  \item Give the zero of $P$. Is your answer an input or an output? What does it mean here?

\writelines{2}

  \item Fill in the table from the graph. Write each interval in interval notation.
\smallskip

\renewcommand{\arraystretch}{1.4}
\begin{tabularx}{\linewidth}{@{}p{2.4cm} p{3.0cm} X@{}}
\rowcolor{mist}
\textbf{Behavior} & \textbf{On these inputs} & \textbf{What it means about the gym} \\
\hline
Constant & \blank{2.2cm} & \blank{6.4cm} \\
Increasing & \blank{2.2cm} & \blank{6.4cm} \\
Decreasing & \blank{2.2cm} & \blank{6.4cm} \\
\end{tabularx}

  \item Give the absolute maximum and the absolute minimum of $P$, each as a \emph{value and a
        location}. Then answer the manager's actual question: \emph{when should the second staff
        member be scheduled?}

\writelines{3}
\end{enumerate}
\end{tcolorbox}

\vspace{0.10in}

\boxguard[30]
\begin{tcolorbox}[colback=white, colframe=black!40,
    title=\textbf{Tier A --- Approaching Proficiency},
    fonttitle=\bfseries, arc=2mm, left=3mm, right=3mm, top=2mm, bottom=2mm, breakable]
\small
Graph (a) shows $M(x)$, a start-up's monthly profit in \emph{thousands} of dollars, $x$ months after
it launched, for its first eight months. Graph (b) shows a function $g$ with no context attached.

\begin{center}
\begin{tabular}{@{}c@{\hspace{0.6cm}}c@{}}
\begin{tikzpicture}
\begin{axis}[
    width=6.8cm, height=4.6cm, title={\scriptsize (a) $M$, a start-up's monthly profit},
    xlabel={$x$ (months)}, ylabel={$M$ (thousands of \$)},
    xmin=0, xmax=8, ymin=-8.4, ymax=10.4,
    xtick={0,1,2,3,4,5,6,7,8}, ytick={-8,-6,-4,-2,0,2,4,6,8,10},
    grid=both, grid style={linegray!45}, axis lines=left,
    tick label style={font=\tiny}, label style={font=\tiny}]
  \addplot[royal, very thick, domain=0:8, samples=90]{-(x-4)^2+9};
\end{axis}
\end{tikzpicture}
&
\begin{tikzpicture}
\begin{axis}[
    width=6.0cm, height=4.6cm, title={\scriptsize (b) $g$},
    xmin=-3.4, xmax=3.4, ymin=0, ymax=9.4,
    xtick={-3,-2,-1,0,1,2,3}, ytick={0,1,2,3,4,5,6,7,8,9},
    grid=both, grid style={linegray!45}, axis lines=left,
    tick label style={font=\tiny}]
  \addplot[linegray, thick, dashed, domain=-3.3:3.3, samples=2]{1};
  \addplot[royal, very thick, domain=-3.3:3, samples=90]{2^x+1};
\end{axis}
\end{tikzpicture}
\end{tabular}
\end{center}

\vspace{-4pt}
{\scriptsize In (b) the dashed gray line is a landmark, not part of the graph of $g$.}

\begin{enumerate}[leftmargin=1.6em, itemsep=4pt, topsep=2pt]
  \item \textbf{Graph (a).} Give the $y$-intercept as a point and the zeros of $M$ as numbers. Then
        say what a zero of $M$ means to the people running this company.

\writelines{3}

  \item \textbf{Graph (a).} Name the interval where $M$ is increasing and the interval where it is
        decreasing. Then give the absolute maximum as a value and a location, and write it as a
        sentence about money and time.

\writelines{3}

  \item \textbf{Graph (a), the careful one.} What is the absolute \emph{minimum} of $M$, and where
        does it happen? Read the whole graph before you answer --- there is more than one right
        location, and you need both.

\writelines{2}

  \item \textbf{Graph (b).} Answer each question, and where the answer is ``none,'' say why in a
        few words.
        \begin{enumerate}[label=(\alph*), leftmargin=1.6em, itemsep=3pt, topsep=3pt]
          \item The $y$-intercept, and the zeros of $g$.

\writelines{2}

          \item The intervals where $g$ is increasing, decreasing, or constant.

\writelines{1}

          \item The absolute maximum and the absolute minimum of $g$.

\writelines{2}
        \end{enumerate}
\end{enumerate}
\end{tcolorbox}

\vspace{0.10in}

\boxguard[30]
\begin{tcolorbox}[colback=white, colframe=black!40, title=\textbf{Tier E --- Extension},
    fonttitle=\bfseries, arc=2mm, left=3mm, right=3mm, top=2mm, bottom=2mm, breakable]
\small

\begin{enumerate}[leftmargin=1.6em, itemsep=4pt, topsep=2pt]
  \item \textbf{Two claims, two wrong axes.} Each classmate below read a real number off a real
        graph and still answered the wrong question. Name the mistake and give the correct answer.
        \begin{enumerate}[label=(\alph*), leftmargin=1.6em, itemsep=3pt, topsep=3pt]
          \item About graph (a) in Tier A: \emph{``The profit is increasing from $-7$ to $9$.''}

\writelines{2}

          \item About the temperature graph in your notes: \emph{``The absolute minimum is $1$
                degree, at $t=12$, because the graph is still going down when it ends.''}

\writelines{2}
        \end{enumerate}

  \item \textbf{Cut the domain and watch three things change.} Suppose the start-up in Tier A graph
        (a) is sold, and the buyer only ever sees months $5$ through $8$ --- that is, the domain is
        cut down to $5 \le x \le 8$ and nothing about the curve is changed.
        \begin{enumerate}[label=(\alph*), leftmargin=1.6em, itemsep=3pt, topsep=3pt]
          \item How many zeros does the buyer see, and how many were there before?

\writelines{1}

          \item Give the absolute maximum and absolute minimum the buyer would report, each with its
                location. Are they at turning points or at endpoints?

\writelines{2}

          \item In one sentence: what would the buyer believe about this company that the founders
                know is false?

\writelines{2}
        \end{enumerate}

  \item \textbf{Justify.} A classmate says: \emph{``Every graph has at least one zero --- it has to
        cross the $x$-axis somewhere.''} Decide whether that is true, and defend your answer with a
        specific graph from today's work. Then give an example of a function whose absolute minimum
        happens at \emph{two} different locations.

\writelines{3}
\end{enumerate}
\end{tcolorbox}

\end{document}
